\subsubsection{0-1 BFS}

\paragraph{Idea intuitiva.}
Variante de BFS para grafos con \textbf{pesos 0 o 1}. Usa un \texttt{deque}: aristas de peso 0 se encolan al frente (\texttt{push\_front}), aristas de peso 1 al final (\texttt{push\_back}). Asi, la primera vez que se extrae un nodo, su distancia es la definitiva, igual que en Dijkstra pero en $O(V + E)$. La intuicion: el deque mantiene los nodos ordenados por distancia; las aristas de peso 0 no cambian la distancia (entran al frente), las de peso 1 la incrementan (entran al final).

\paragraph{Propiedades clave (invariantes).}
\begin{itemize}\setlength\itemsep{0pt}
	\item Asume pesos \textbf{solo} 0 o 1.
	\item Mejor caso que Dijkstra (sin factor $\log V$).
	\item No funciona con pesos $\ge 2$.
	\item La cola mantiene los nodos ordenados por distancia no decreciente.
\end{itemize}

\paragraph{Codigo clave.}
El deque doble:
\begin{verbatim}
deque<int> q;
q.push_front(source);
dist[source] = 0;
while (!q.empty()) {
    int v = q.front(); q.pop_front();
    for (auto [w, u] : adj[v]) {
        if (dist[v] + w < dist[u]) {
            dist[u] = dist[v] + w;
            if (w == 0) q.push_front(u);
            else q.push_back(u);
        }
    }
}
\end{verbatim}

\paragraph{Complejidad.}
\begin{itemize}\setlength\itemsep{0pt}
	\item $O(V + E)$.
	\item Cada nodo se inserta a lo sumo $O(1)$ veces (con check de distancia).
\end{itemize}

\paragraph{Donde tocar para modificar.}
\begin{itemize}\setlength\itemsep{0pt}
	\item \textbf{Pesos en $[0, k]$ small}: usar small k-ary BFS o Dijkstra.
	\item \textbf{Reconstruir camino}: aniadir \texttt{parent[]}.
	\item \textbf{Pesos negativos}: no funciona, usar Bellman-Ford.
\end{itemize}

\paragraph{Trampas y errores frecuentes.}
\begin{itemize}\setlength\itemsep{0pt}
	\item Si los pesos son otro rango, el algoritmo da respuestas incorrectas. Verificar que las aristas son solo 0/1.
	\item El chequeo \texttt{dist[v] + w < dist[u]} es crucial; sino, el algoritmo procesa nodos multiples veces.
	\item Si una arista tiene peso $\ge 2$, aniadir verificaciones o cambiar a Dijkstra.
\end{itemize}

\paragraph{Codigo.}
Ver \texttt{cpp.tex}, seccion \textit{0-1 BFS}.
